% ---------------------------------------------------------------------------
% Response B, anatomized: the relation graph and the MRF it induces, side by
% side, with the four LCS readouts and the per-claim marginals below them.
%
% The point of putting the two views next to each other is that they are the
% SAME object at two levels of resolution: every edge on the left becomes one
% black pairwise factor on the right, and the five priors that are implicit on
% the left become the five grey unary factors that are explicit on the right.
% The score panel then shows what inference over that MRF returns.
%
% Both panels are drawn at one shared scale (\rbx/\rby) rather than each in its
% own \resizebox, so neither can appear larger than the other. Node, edge and
% relw=p styles live in preamble.tex, shared with figures/graph-coherence-pair,
% mrf-coherence-pair and motivating-example, so no panel can drift.
%
% Every number is produced by scripts/lcs_worked_pair.py from
% data/lcs/example-7-coherence-pair.json -- exact, 2^5 = 32-world enumeration,
% computed twice by independent paths and asserted equal.
% ---------------------------------------------------------------------------
\begin{figure}[tbp]
\centering
\newcommand{\rbx}{13.4mm}
\newcommand{\rby}{11.8mm}

%% ===== (a) the relation graph | (b) the MRF ==============================
\begin{tabularx}{\textwidth}{@{}X@{\hspace{4pt}}X@{}}
\begin{minipage}[b]{\linewidth}\centering
\begin{tikzpicture}[x=\rbx, y=\rby]
  \node[atomtrue]  (a1) at (0,0)        {$a_1$};
  \node[atomtrue]  (a2) at (1.2,0)      {$a_2$};
  \node[atomtrue]  (a3) at (2.4,0)      {$a_3$};
  \node[atomfalse] (a4) at (0.6,-1.6)   {$a_4$};
  \node[atomfalse] (a5) at (1.85,-1.6)  {$a_5$};

  \node[prlab,above=0.4mm of a1] {$\pi{=}.9$};
  \node[prlab,below=0.4mm of a2] {$\pi{=}.9$};
  \node[prlab,above=0.4mm of a3] {$\pi{=}.9$};
  \node[prlab,left=0.5mm of a4]  {$\pi{=}.1$};
  \node[prlab,right=0.5mm of a5] {$\pi{=}.1$};

  \draw[ent,relw=0.90] (a1) -- (a2) node[wlab,above]{$.90$};
  \draw[con,relw=0.90] (a2) -- (a3) node[wlab,above]{$.90$};
  \draw[ent,relw=0.88] (a4) -- (a3)
        node[wlab,above left=-0.4mm and -1.1mm,pos=0.32]{$.88$};
  \draw[con,relw=0.90] (a1) -- (a4) node[wlab,left,pos=0.55]{$.90$};
  \draw[ent,relw=0.85] (a4) -- (a5) node[wlab,below]{$.85$};
  \draw[con,relw=0.88] (a2) to[out=-8,in=25]
        node[wlab,right,pos=0.55,inner sep=1pt]{$.88$} (a5);

  \path (0.6,-2.35) -- (0.6,0.95);   % strut: match the right panel's depth
\end{tikzpicture}\\[2pt]
{\scriptsize \textbf{(a) The relation graph:} three \ent{} (solid blue),\\
 three \con{} (dashed red); width grows with $p$.}
\end{minipage}
&
\begin{minipage}[b]{\linewidth}\centering
\begin{tikzpicture}[x=\rbx, y=\rby]
  \node[atomtrue]  (a1) at (0,0)        {$a_1$};
  \node[atomtrue]  (a2) at (1.15,0)     {$a_2$};
  \node[atomtrue]  (a3) at (2.3,0)      {$a_3$};
  \node[atomfalse] (a4) at (1.15,-1.3)  {$a_4$};
  \node[atomfalse] (a5) at (2.3,-1.3)   {$a_5$};

  % unary factors phi_i = [1-pi_i, pi_i]: the priors, made explicit
  \node[unary] (p1) at (0,0.78)     {}; \draw[flink] (p1) -- (a1);
  \node[unary] (p2) at (1.15,0.78)  {}; \draw[flink] (p2) -- (a2);
  \node[unary] (p3) at (2.3,0.78)   {}; \draw[flink] (p3) -- (a3);
  \node[unary] (p4) at (1.15,-2.08) {}; \draw[flink] (p4) -- (a4);
  \node[unary] (p5) at (2.3,-2.08)  {}; \draw[flink] (p5) -- (a5);
  \node[prlab,left=0.4mm of p1] {$\phi_1$};
  \node[prlab,left=0.4mm of p2] {$\phi_2$};
  \node[prlab,left=0.4mm of p3] {$\phi_3$};
  \node[prlab,left=0.4mm of p4] {$\phi_4$};
  \node[prlab,left=0.4mm of p5] {$\phi_5$};

  % Pairwise factors psi_r, one per mined relation. The two crossing links
  % (a4->a3 and a2->a5) are routed around opposite sides so neither factor
  % sits on the other's edge.
  \node[factor] (f12) at (0.575,0)      {}; \draw[flink] (a1) -- (f12) -- (a2);
  \node[factor] (f23) at (1.725,0)      {}; \draw[flink] (a2) -- (f23) -- (a3);
  \node[factor] (f43) at (1.72,-0.62)   {}; \draw[flink] (a4) -- (f43) -- (a3);
  \node[factor] (f14) at (0.42,-0.80)   {}; \draw[flink] (a1) -- (f14) -- (a4);
  \node[factor] (f45) at (1.725,-1.3)   {}; \draw[flink] (a4) -- (f45) -- (a5);
  \node[factor] (f25) at (2.82,-0.62)   {};
  \draw[flink] (a2) to[out=-32,in=118] (f25);
  \draw[flink] (f25) to[out=-118,in=32] (a5);
  \node[prlab,above=0.3mm of f12] {$\psi_1$};
  \node[prlab,above=0.3mm of f23] {$\psi_2$};
  \node[prlab,left=0.4mm of f43]  {$\psi_3$};
  \node[prlab,left=0.3mm of f14]  {$\psi_4$};
  \node[prlab,below=0.3mm of f45] {$\psi_5$};
  \node[prlab,right=0.4mm of f25] {$\psi_6$};
\end{tikzpicture}\\[2pt]
{\scriptsize \textbf{(b) The MRF it induces:} one $\psi_r$ per edge (black),\\
 one $\phi_i$ per prior (grey).}
\end{minipage}
\end{tabularx}

\medskip
%% ===== (c) the scores ====================================================
% Bands (c) and the per-claim gloss are paper-only: the standalone PNG built by
% scripts/make_response_b_png.sh keeps just the two plots, so \rbnocaption gates
% this block and the \caption below alike.
\ifdefined\rbnocaption\else
\begin{tabularx}{\textwidth}{@{}l@{\hspace{7pt}}c@{\hspace{7pt}}X@{}}
\toprule
Readout & Value & What it says \\
\midrule
$\LCS_{\mm}$ (selected) & $\mathbf{0.494}$
  & Mean posterior marginal: the expected fraction of B's claims its own argument
    upholds. Against a prior-only mean of $0.580$, the arrangement \emph{costs}
    B support rather than adding any. \\
$\LCS_{\cons}$ & $0.414$
  & Two-term consistency: little over a coin flip that no asserted conflict is
    realized, against $0.963$ for the coherent arrangement of the same claims. \\
$\LCS_{\rei}$ & $0.348$
  & The reified-node variant, for comparison; it agrees on the ordering. \\
$\log Z$ & $-3.78$
  & Global mass surviving the constraints ($-0.95$ for the coherent
    arrangement): the six factors are jointly hard to satisfy. \\
\midrule
$\LCS_{\lp}$ & $(0.267)$
  & Parenthesized because it is \emph{not} comparable across responses: its
    references $\log Z_{\max}=-1.65$ and $\log Z_{\min}=-4.55$ are rebuilt from
    B's own edge skeleton, so it reports distance to \emph{that} ceiling. A
    within-response diagnostic only. \\
\bottomrule
\end{tabularx}

\smallskip
{\footnotesize\raggedright
\textbf{Per-claim marginals} (the diagnostic the scalar cannot give):
$a_1\,0.960$, $a_2\,0.973$, $\mathbf{a_3\,0.513}$, $a_4\,0.006$, $a_5\,0.018$.
Three of five finish below their own prior. The one that matters is $a_3$: a claim
\emph{true of the world}, with prior $0.9$, dragged to a coin flip. Its only support
is $\psi_3$, arriving from the false $a_4$ --- which $\psi_4$ has the response
denying. A premise the text refuses cannot transmit support, so $a_3$ is left
holding an argument that undercuts itself. All values exact, by $2^5{=}32$-world
enumeration.\par}

\caption{\textbf{Response B, anatomized.} The incoherent arrangement of the
  five-claim pair (\S\ref{sec-example-pair}) at two levels of resolution, and what
  inference over it returns. Left: the relation graph, blue nodes true of the world
  and red false. Right: the Markov random field of Def.~\ref{def-mrf} that graph
  induces, as an explicit factor graph --- the mapping is one-to-one, six edges to
  six pairwise factors $\psi_r$, five priors to five unary factors $\phi_i$. Below:
  the four readouts of \S\ref{sec-readouts} and the per-claim marginals. The reading to
  take from the pair is that the $\phi_i$ are identical to those of the coherent
  arrangement --- the claim set and every factuality verdict are unchanged --- so the
  whole of the score difference is carried by the $\psi_r$, and the whole of the
  damage is localized by the marginals to $a_3$.}
\label{fig:response-b-anatomy}
\fi
\end{figure}
